Для основной оценки использован графический метод п.~9 описания [1, с.~3]
с заменой построений линейным МНК. По точкам $U\ge16$~В и $U\le-16$~В
найдены две асимптоты $L_+(U)$ и $L_-(U)$, а по $|U|\le4{,}6$~В ---
общая центральная прямая $C(U)=gU+b$. Абсциссу центра $U_0$ определили
из условия $C(U_0)=[L_+(U_0)+L_-(U_0)]/2$.
На графиках $U'=U-U_0$, $I'=I-C(U_0)$,
а $\Is=[L_+(U_0)-L_-(U_0)]/2$.
Вычисленное по \eqref{eq:temp} значение соответствует половине расстояния
по горизонтали от оси $U'=0$ до пересечения касательной с уровнем $I'=\Is$.

Стандартные ошибки трёх прямых оценены по остаткам МНК. При переносе
в вычисляемые величины использована ковариация наклона и свободного члена
каждой прямой. Центральные и крайние точки принадлежат разным выборкам.
Влияние разрыва при смене полярности проверено повторной обработкой
с отдельным свободным членом для каждого положения П2.
